% =====================================================================
%  Chapter 5 - RQ2. Replaces the "RQ2 Planning" chapter entirely.
%  Structure follows Chapter 4: Aim, Implementation, Experimental Setup,
%  Results, Analysis, Answer. Results is divided three ways because RQ2
%  has three sub-questions and the experiment has three phases.
%
%  Values marked [PENDING] come from Phase 2 and are not yet collected.
%  Everything else is measured, from the five baseline sessions.
% =====================================================================

\section{Autonomous Threshold-Based Fault Detection on a Physical Remote
Autonomous System (RQ2)}\label{sec:rq2}

\subsection{Aim}\label{sec:rq2-aim}

RQ2 asks: \emph{How can system metrics such as battery level, connectivity and
performance be used to detect faults in remote autonomous systems?} Study 1
established that the pipeline can observe health state, but the health state it
observed was supplied by a dataset label rather than computed from the
telemetry. This chapter removes that dependency. The same pipeline is connected
to a physical DJI Tello, each threshold is derived from the aircraft's own
measured baseline, and Prometheus alerting rules determine health state from raw
readings with no labels available to them.

This section documents the instrument built for that purpose, the three-phase
experiment carried out with it, and the results measured against the criteria
fixed in Section~\ref{tab:eval-rq2}. The phases correspond to the three
sub-questions. Phase 1 derives the thresholds and answers RQ2.1, Phase 2 tests
whether the rules detect induced faults and answers RQ2.2, and Phase 3
quantifies detection latency and false positive rate for RQ2.3.

% =====================================================================
\subsection{Implementation}\label{sec:rq2-impl}

Study 2 required three capabilities that Study 1 did not: repeatable physical
flights, thresholds computed from measured data, and alert timings that can be
recovered after a run has finished. Each is provided by a separate tool, and all
are published in the project repository.

\subsubsection{Session Harness}\label{sec:rq2-harness}

\texttt{flight\_session.py} is the instrument used for every physical flight. A
single process performs three jobs, because the Tello accepts one SDK client at
a time and a second client competing for the same UDP ports would prevent both
from working.

\begin{enumerate}
    \item It flies a fixed manoeuvre script consisting of take-off, hover,
    translation, rotation, hover and land, so that variation between sessions
    reflects the platform rather than the pilot.
    \item It records raw telemetry to a CSV file at \SI{1}{\hertz} and tags
    every sample with the manoeuvre step that was executing when the sample was
    taken.
    \item It publishes the same readings to Prometheus on port 8001.
\end{enumerate}

Manual piloting was considered and rejected. Three repetitions of a fault must
be the same fault for the repetitions to be meaningful, and hand-flown runs
would introduce operator variation into every measurement. The single-client
constraint also rules it out in practice, since flying from the mobile
application would occupy the connection the exporter requires.

The harness writes a separate event log containing manoeuvre boundaries, stop
reasons and operator markers. Pressing \textsc{Enter} during a flight records
$t_{\text{induce}}$. This is necessary for the link degradation runs, because
only the operator knows the moment at which they began increasing separation.

Four modes are provided. The \texttt{baseline} mode flies the Phase 1 script.
The \texttt{battery} mode disables the harness's own landing so that the
aircraft's low-battery failsafe ends the run, since that failsafe is the object
of measurement rather than a constraint to be avoided. The \texttt{thermal} mode
flies continuously at \SI{100}{\centi\metre\per\second} with no hover segment.
The \texttt{link} mode holds a stationary hover while the operator moves.

Two pre-flight checks refuse to begin a session that could not be compared with
the others. The first rejects a battery below \SI{80}{\percent} and the second,
which is skipped for thermal runs, rejects an onboard temperature above
\SI{60}{\celsius}. Both were added after sessions were lost to those conditions,
and the reasoning behind the second is given in Section~\ref{sec:rq2-thresholds}.

\subsubsection{Metric Set and Platform Constraints}\label{sec:rq2-metrics}

The fields the Tello exposes are listed in Table~\ref{tab:tello-fields}.
Baseline collection established that several of them behave differently from
what their SDK field names imply, and those findings are reported in
Section~\ref{sec:rq2-platform}. The metric set carried through to the alerting
rules is battery percentage, onboard temperature, attitude, barometric height
and target liveness.

\subsubsection{Alerting Rules}\label{sec:rq2-rules}

Fault detection is performed entirely by Prometheus alerting rules evaluated
against the raw metrics. No metric published by the exporter is a health label,
and there is no equivalent of the \texttt{ras\_health\_status} metric used in
Study 1. Every rule carries a \texttt{for: 5s} hold, so that a single sample
crossing a threshold cannot raise an alert. The rules are presented in
Section~\ref{sec:rq2-thresholds} together with the values derived for them.

\subsubsection{Analysis Tooling}\label{sec:rq2-tooling}

Three further scripts make the results reproducible from the archived logs.

\begin{itemize}
    \item \texttt{check\_session.py} validates a session immediately after
    landing. It checks sample coverage against the expected \SI{1}{\hertz},
    gaps between samples, battery monotonicity, SNR coverage, manoeuvre
    completeness and metadata. A session that fails is re-flown rather than
    retained.
    \item \texttt{derive\_thresholds.py} pools the baseline sessions, applies
    the derivation rules given in Table~\ref{tab:threshold-rules}, and writes
    both the Prometheus rules file and a machine-readable copy of the same
    values.
    \item \texttt{collect\_alerts.py} recovers alert firing times from the
    \texttt{ALERTS} series that Prometheus maintains for its own alert state,
    computes threshold-crossing times from the telemetry log, and reports
    detection latency. Both calculations read the machine-readable threshold
    file, so the onset calculation and the alerting rules cannot diverge.
\end{itemize}

% =====================================================================
\subsection{Experimental Setup}\label{sec:rq2-setup}

\subsubsection{Platform and Environment}\label{sec:rq2-platform-setup}

All sessions used a single DJI Tello, serial number 0TQZGAJED007GD, flown
indoors at a private residence at an ambient temperature of \SI{21}{\celsius}.
One airframe was used throughout so that the thresholds and the fault runs refer
to the same unit. The consequences of that decision are addressed in
Section~\ref{subsec:method-threats}.

Two batteries, designated B1 and B2, were alternated between sessions so that
battery identity would not coincide with session order. Prometheus, Grafana, the
harness and the exporter all ran on a single host, so every timestamp comes from
one system clock and clock synchronisation between devices cannot introduce
measurement error. The scrape interval and the rule evaluation interval were
both set to \SI{1}{\second}.

\subsubsection{Session Protocol}\label{sec:rq2-protocol}

Each session ran for four minutes and began on a fully charged battery, with the
aircraft returned to ambient temperature beforehand. Four-minute sessions sample
only the upper region of the discharge curve. The lower region, approaching the
automatic low-battery landing, is characterised instead by the Phase 2 depletion
runs, which continue until that failsafe engages.

Lighting was controlled alongside ambient temperature. The Tello holds position
using a downward-facing camera and a time-of-flight sensor, and a session
attempted in a dimly lit room produced erratic readings and uncommanded drift.
That result is reported in Section~\ref{subsec:method-threats}.

\subsubsection{Experimental Runs}\label{sec:rq2-runs}

Table~\ref{tab:rq2-runs} lists every run in the study. Sessions excluded from
analysis are retained in the repository with the reason for exclusion recorded,
following the reliability commitment stated in Section~\ref{subsec:method-threats}.

\begin{table}[htbp]
\centering
\caption{Study 2 Experimental Runs}
\label{tab:rq2-runs}
\small
\begin{tabular}{|l|l|l|p{5.0cm}|l|}
\hline
\textbf{Phase} & \textbf{Runs} & \textbf{Mode} & \textbf{Purpose} & \textbf{Sub-question} \\
\hline
1 & BASE-01--05 & baseline & Characterise normal operation and derive thresholds & RQ2.1 \\
\hline
-- & FP-01--03 & baseline & Held-out normal flights for the false positive rate & RQ2.3 \\
\hline
2 & FAULT-BAT-01--03 & battery & Flight continues until the low-battery failsafe engages & RQ2.2, RQ2.3 \\
\hline
2 & FAULT-LINK-01--03 & link & Operator increases separation and obstructs line of sight & RQ2.2, RQ2.3 \\
\hline
2 & FAULT-THERM-01--03 & thermal & Consecutive sorties flown without a cool-down interval & RQ2.2, RQ2.3 \\
\hline
\end{tabular}
\end{table}

% =====================================================================
\subsection{Phase 1 Results: Baseline Characterisation
(RQ2.1)}\label{sec:rq2-phase1}

\subsubsection{Session Summary}\label{sec:rq2-sessions}

Five baseline sessions were flown and all five passed validation, producing 1274
pooled in-flight samples. Table~\ref{tab:baseline-sessions} summarises them.
Sample coverage was \SI{100}{\percent} in every session and no gap between
consecutive samples exceeded \SI{1.1}{\second}, so the distributions reported
below are computed from complete data.

\begin{table}[htbp]
\centering
\caption{Phase 1 Baseline Sessions}
\label{tab:baseline-sessions}
\small
\begin{tabular}{|l|c|c|c|c|c|c|}
\hline
\textbf{Session} & \textbf{Battery} & \textbf{Ambient} & \textbf{Duration} &
\textbf{Battery span} & \textbf{Depletion} & \textbf{In-flight temp} \\
 & & (\si{\celsius}) & (s) & (\%) & (\%/min) & (\si{\celsius}) \\
\hline
BASE-01 & B1 & 21 & 271 & 100--66 & 7.53 & 56--63 \\
BASE-02 & B2 & 21 & 270 & 88--40 & 10.67 & 57--62 \\
BASE-03 & B1 & 21 & 273 & 100--66 & 7.47 & 56--62 \\
BASE-04 & B2 & 21 & 273 & 92--45 & 10.33 & 61--65 \\
BASE-05 & B1 & 21 & 275 & 100--69 & 6.76 & 58--63 \\
\hline
\end{tabular}
\end{table}

\subsubsection{Platform Behaviour}\label{sec:rq2-platform}

Baseline collection produced six findings about the platform, each of which
changes how the metric concerned can be used. They are reported here rather than
in the analysis because they form part of the answer to RQ2.1: establishing
which metrics represent operational health requires knowing how those metrics
actually behave, and on this platform several of them do not behave as their
field descriptions suggest.

\paragraph{Onboard temperature settles toward an operating point.} Temperature
does not vary about a fixed mean during a session. It converges toward a steady
operating value of approximately \SI{62}{\celsius}, and the direction of change
depends on where the aircraft starts. Sessions beginning at ambient warmed from
\SI{55}{}--\SI{58}{\celsius} to \SI{62}{\celsius} over four minutes. Session
BASE-04, which began at \SI{65}{\celsius} because it followed an earlier flight,
cooled to the same point instead. Table~\ref{tab:temp-settling} shows the
progression.

\begin{table}[htbp]
\centering
\caption{Onboard Temperature by Elapsed Flight Time (\si{\celsius})}
\label{tab:temp-settling}
\small
\begin{tabular}{|l|c|c|c|c|c|c|}
\hline
\textbf{Session} & \textbf{Start} & \SI{0.5}{min} & \SI{1}{min} &
\SI{2}{min} & \SI{3}{min} & \SI{4}{min} \\
\hline
BASE-01 & 55 & 58.8 & 60.3 & 61.8 & 62.2 & 62.5 \\
BASE-02 & 58 & 58.1 & 59.5 & 60.2 & 61.3 & 61.5 \\
BASE-03 & 56 & 58.2 & 59.5 & 60.7 & 61.2 & 61.9 \\
BASE-04 & 65 & 62.3 & 62.1 & 62.5 & 62.9 & 62.8 \\
BASE-05 & 58 & 60.2 & 60.7 & 61.9 & 62.1 & 62.2 \\
\hline
\end{tabular}
\end{table}

Two consequences follow. The first concerns the assumed threshold carried into
Study 2. A ceiling of \SI{60}{\celsius}, taken from industrial motor-monitoring
practice, sits below the operating point the aircraft reaches in normal flight,
so the rule fires on every healthy session once the airframe has warmed. That
is the false positive observed during preliminary testing.

The second concerns the derivation rule, and is treated in
Section~\ref{sec:rq2-thresholds}. Because a four-minute session spends most of
its samples in the warm-up transient rather than at the operating point, a
percentile taken over the whole session does not describe steady operation.

\paragraph{Barometric readings track atmospheric pressure rather than altitude.}
Absolute barometer values moved by approximately \SI{20}{\centi\metre} between
flight days and by \SI{3.6}{\centi\metre} between sessions flown in a single
room. The aircraft's actual ground clearance was unchanged over the same period.
A threshold expressed as an absolute band would therefore encode the prevailing
weather rather than any property of the aircraft. The metric is usable only as
within-session drift, and Section~\ref{sec:rq2-thresholds} derives it that way.

\paragraph{Reported height is a derived quantity, not a measured one.} The
median value returned by \texttt{get\_height()} fell from \SI{60}{\centi\metre}
to \SI{30}{\centi\metre} across successive sessions and reached
\SI{-40}{\centi\metre}, while time-of-flight ground clearance held between
\SI{77}{} and \SI{78}{\centi\metre} throughout. The aircraft maintained
altitude and the sensor drifted. Height and barometer are therefore treated as
drift indicators and time-of-flight as the altitude signal, which reverses the
roles assigned to them in Table~\ref{tab:tello-fields}. The reassignment follows
from measured behaviour rather than from the SDK documentation.

\paragraph{Battery identity dominates the depletion rate.} Battery B1 depleted
at a mean rate of \SI{7.25}{\percent\per\minute} across three sessions, with
individual runs at 7.53, 7.47 and \SI{6.76}{\percent\per\minute}. Battery B2
depleted at \SI{10.50}{\percent\per\minute} across two sessions, at 10.67 and
\SI{10.33}{\percent\per\minute}. The two distributions do not overlap. B2 also
never charged above \SI{92}{\percent} while B1 reached \SI{100}{\percent} in
every session, which gives a second and independent indication of a degraded
cell. A single depletion threshold cannot describe both batteries, so the
derivation in Section~\ref{sec:rq2-thresholds} takes the worse of the two.
Alternating batteries between sessions is what made the difference visible.

\paragraph{Wi-Fi SNR carries no usable variance.} The metric returned a constant
value of 90 across all 1274 in-flight samples. A percentile of a constant is
that constant, so the derivation rule for bounded metrics degenerates and no
threshold can be established. SNR is additionally exposed as a command-channel
query rather than as part of the broadcast state stream, so it cannot be sampled
at \SI{1}{\hertz} without contending with flight commands. It was polled at
\SI{5}{\second} intervals and carried forward between polls, with the staleness
of each carried value recorded alongside it.

\paragraph{Velocity fields are uninformative at the sampling rate used.}
Reported velocity remained close to zero even during commanded translation,
because each movement completes between samples. The manoeuvre is instead
visible in attitude, where pitch reaches \SI{8}{\degree} during translation
compared with \SI{2}{\degree} in hover.

\subsubsection{Derived Thresholds}\label{sec:rq2-thresholds}

Table~\ref{tab:derived-thresholds} gives the threshold derived for each
monitored metric, together with the metric class it belongs to in
Table~\ref{tab:threshold-rules} and the derivation that produced it. Every value
follows from a stated rule applied to the pooled baseline distribution and none
is assumed.

\begin{table}[htbp]
\centering
\caption{Thresholds Derived from Phase 1 Baseline Data}
\label{tab:derived-thresholds}
\footnotesize
\begin{tabular}{|p{2.4cm}|p{2.2cm}|p{2.0cm}|p{6.4cm}|}
\hline
\textbf{Alert} & \textbf{Metric class} & \textbf{Value} & \textbf{Derivation} \\
\hline
DroneLowBattery & Monotonic & \SI{31}{\percent} &
Firmware auto-land point of \SI{10}{\percent} plus two minutes of lead time at
the worst observed depletion rate of \SI{10.50}{\percent\per\minute} \\
\hline
DroneOverTemperature & Settling & \SI{64}{\celsius} &
99th percentile of the settled region, taken as the final \SI{90}{\second} of
each session, applied in the high direction only \\
\hline
DroneUnstableAttitude & Variance & 3.87 pitch, 1.20 roll &
99th percentile of the \SI{5}{\second} rolling standard deviation \\
\hline
DroneBarometricDrift & Drift & \SI{0.31}{\centi\metre} &
99th percentile of the deviation from the metric's own \SI{60}{\second} rolling
mean \\
\hline
DronePollFailure & Liveness & n/a &
Binary. Missing data is itself the fault condition and no value is derived \\
\hline
\end{tabular}
\end{table}

Two of the rules required guards that the derivation itself made necessary, and
both are stated here because they affect what the thresholds mean.

The temperature rule is evaluated only while the aircraft is airborne, using the
condition \texttt{ras\_tof\_cm > 20}. The thresholds are derived from in-flight
samples, but \SI{12}{\percent} of ground samples in the baseline exceed the
resulting ceiling on a healthy aircraft, for the reason given in
Section~\ref{sec:rq2-platform}. Without the guard the rule fires before every
take-off.

The same rule is applied in one direction only. A two-sided rule using the 1st
percentile of \SI{57}{\celsius} fires during the first seconds of healthy flight,
while the aircraft is still below its operating temperature. Session BASE-01
recorded \SI{56}{\celsius} for six consecutive seconds in that period, which
exceeds the \SI{5}{\second} hold and would therefore have raised an alert. Since
no mechanism on this platform produces an under-temperature fault, a lower bound
contributes false positives without contributing any detection capability.
Table~\ref{tab:threshold-rules} permits this by specifying the percentile per
metric direction.

The threshold is derived from the settled region rather than from the whole
session, for the reason established in Section~\ref{sec:rq2-platform}. The
distinction is not incidental. A percentile taken across all in-flight samples
gives \SI{63}{\celsius}, which lies inside the \SI{62}{}--\SI{64}{\celsius}
band the aircraft occupies once settled, because the warm-up transient
contributes the majority of samples in a four-minute session and pulls the
percentile down. That value was applied and a held-out normal flight crossed it,
raising an alert on a healthy aircraft. Restricting the derivation to the final
\SI{90}{\second} of each session gives \SI{64}{\celsius}, which no baseline
session and no other held-out flight exceeds.

The reasoning is the same one Table~\ref{tab:threshold-rules} already applies to
battery level. A quantity that trends through a session by design is not
described by a percentile of its distribution, because that distribution is
dominated by the trend rather than by the operating state. Battery level trends
downward and temperature trends toward an operating point, and neither is
stationary in the sense the bounded rule assumes.

The barometric drift rule is evaluated only once its entire \SI{60}{\second}
window consists of airborne samples, using the condition
\texttt{min\_over\_time(ras\_tof\_cm[60s]) > 20}. While the window is still
filling after take-off its mean remains dominated by ground-level readings, and
the deviation therefore measures that transition rather than the aircraft. The
value reads \SI{0.68}{\centi\metre} immediately after take-off and settles below
\SI{0.31}{\centi\metre} once the window is complete.

\subsubsection{Assumed and Derived Thresholds Compared}\label{sec:rq2-assumed}

Table~\ref{tab:assumed-vs-derived} sets the thresholds assumed before baseline
collection against those derived from it. Both configurations are published in
the repository so that the comparison can be verified.

\begin{table}[htbp]
\centering
\caption{Assumed and Derived Thresholds Compared}
\label{tab:assumed-vs-derived}
\small
\begin{tabular}{|p{2.6cm}|p{3.2cm}|p{3.4cm}|p{4.0cm}|}
\hline
\textbf{Condition} & \textbf{Assumed} & \textbf{Derived} & \textbf{Consequence of the assumed value} \\
\hline
Low battery & $<\SI{15}{\percent}$ & $<\SI{31}{\percent}$ &
Leaves under \SI{30}{\second} of lead time before the failsafe at the observed
depletion rate \\
\hline
Temperature & $>\SI{60}{\celsius}$ & $>\SI{63}{\celsius}$, airborne only &
Fires during the healthy first seconds of every flight \\
\hline
Attitude & $|\text{pitch}|$ or $|\text{roll}| > \SI{30}{\degree}$ &
\SI{5}{\second} rolling standard deviation $>$ 3.87 or 1.20 &
Baseline attitude never exceeds \SI{8}{\degree}, so the rule cannot fire and
detects nothing \\
\hline
\end{tabular}
\end{table}

The two failures are not equivalent, and the difference between them is worth
stating. The assumed temperature rule produced a visible false positive and was
noticed during preliminary testing. The assumed attitude rule produced no alerts
at all, which is indistinguishable from correct operation until a fault occurs
and is not detected. An assumed threshold can therefore fail in either
direction, and only one of the two failure modes is self-announcing.

% =====================================================================
\subsection{Phase 2 Results: Fault Detection (RQ2.2)}\label{sec:rq2-phase2}

% [PENDING] Complete after the nine fault runs.

\subsubsection{Held-Out Normal Flights}\label{sec:rq2-fp}

% [PENDING] FP-01 to FP-03. Report any alerts raised and the total flight time.

\subsubsection{Battery Depletion}\label{sec:rq2-fault-battery}

% [PENDING] Per run: t_onset, t_alert, and the lead time before the failsafe
% engaged.

\subsubsection{Link Degradation}\label{sec:rq2-fault-link}

% [PENDING] Per run: t_induce, the behaviour of SNR, and whether
% DronePollFailure was raised.

\subsubsection{Thermal and Load Stress}\label{sec:rq2-fault-thermal}

% [PENDING] Per run: starting temperature, in-flight range, alerts raised.

\subsubsection{Detection Outcomes}\label{sec:rq2-confusion}

% [PENDING] Confusion matrix across the nine runs, classifying each as a true
% positive, a false negative or a misclassification. Detection rate against the
% target of 8 of 9.

% =====================================================================
\subsection{Phase 3 Results: Evaluation (RQ2.3)}\label{sec:rq2-phase3}

\subsubsection{Detection Latency}\label{sec:rq2-latency}

% [PENDING] Median, interquartile range and maximum of t_alert minus t_onset
% across all detected faults. State the two lower bounds alongside the result:
% the 1 s scrape interval and the 5 s hold.

\subsubsection{False Positive Rate}\label{sec:rq2-fp-rate}

% [PENDING] False alerts per 10 minutes of normal flight, measured on the
% held-out flights and on the pre-induction segment of each fault run.

\subsubsection{Metric Coverage}\label{sec:rq2-coverage}

% [PENDING] Proportion of monitored metrics present and non-stale, measured with
% up and absent(). The baseline sessions achieved 100 percent sample coverage.

\subsubsection{Evaluation Criteria Results}\label{sec:rq2-criteria}

% [PENDING] Follow the format of Table 12 in Chapter 4.

% =====================================================================
\subsection{Analysis}\label{sec:rq2-analysis}

The claim that Study 2 is built on is that thresholds must be derived from the
platform rather than assumed, and the Phase 1 results support that claim
independently of how the fault runs turn out. Of the three assumed thresholds
carried into Study 2, one fired continuously during healthy operation and one
could not fire at all. Neither was chosen carelessly. Both were conventional
values for the quantity being measured, and what made them unsuitable was the
platform they were applied to. The only way that became apparent was by
measuring it.

The platform findings reported in Section~\ref{sec:rq2-platform} extend the same
point. Three of the fields listed in Table~\ref{tab:tello-fields} do not behave
as their names imply. Temperature converges toward an operating point rather
than varying about a mean, reported height is a derived value that drifts while
true altitude is held, and velocity carries no information at a \SI{1}{\hertz}
sampling rate. A study conducted entirely on replayed data would have inherited
all three errors, because a dataset cannot contradict the description of its own
fields. This is the practical difference between Study 1 and Study 2, and it is
a difference in what the evidence is capable of showing rather than only in the
platform used.

A further observation concerns the derivation procedure itself rather than any
individual threshold. Two metrics were initially placed in the wrong class in
Table~\ref{tab:threshold-rules}, and in both cases the error was invisible until
a threshold derived under the wrong rule was applied to a held-out flight.
Barometric height was treated as bounded, but its absolute value tracks
atmospheric pressure, so the derived band encoded the prevailing weather.
Temperature was also treated as bounded, but it trends toward an operating point,
so the derived ceiling sat below the range of normal steady operation. Both
produced alerts on healthy aircraft, and neither was detectable from inspecting
the derivation rule in isolation.

This suggests that assigning a metric to the correct behavioural class is a
larger source of error than the choice of percentile within a class. The
literature reviewed in Section~\ref{sec:lit-thresholds} concentrates on
threshold values and on the methods used to compute them, and gives less
attention to the prior question of whether the quantity being thresholded is
stationary at all. On this platform that question decided the outcome twice.
Determining it required measuring the platform under normal operation, holding
out further normal flights, and testing the derived thresholds against data that
had taken no part in producing them.

\subsubsection{Sensitivity to Ambient Temperature}\label{sec:rq2-ambient}

The temperature threshold is expressed as an absolute value, which makes it
sensitive to the ambient conditions under which the baseline was collected. All
five baseline sessions were flown at an ambient temperature of
\SI{21}{\celsius}, and their settled means fall between \SI{61.4}{} and
\SI{62.8}{\celsius}. A held-out normal flight conducted in a room estimated to
be \SI{1}{}--\SI{2}{\celsius} warmer settled at \SI{64.0}{\celsius}, above
every baseline session and above the derived ceiling.

The aircraft was in normal condition for that flight. It had cooled overnight,
began on a fully charged battery, and its battery depletion rate
(\SI{7.08}{\percent\per\minute}) and mean acceleration magnitude were within the
baseline range, so there is no indication that it was working harder than in any
other session. The most consistent explanation is that the ambient temperature
displaces the whole operating band, approximately one degree for one degree.

This is a limitation of absolute thermal thresholds rather than of the
derivation procedure, and it constrains the conditions under which the derived
value transfers. The methodology anticipates the mechanism, since ambient
temperature is recorded per session precisely because it shifts the thermal
baseline, but a baseline collected at a single ambient temperature cannot
quantify the relationship. Characterising it would require baseline sessions
across a range of ambient temperatures, or a rule expressed relative to a
measured ambient reference, which the platform does not report. Both are
identified as future work in Section~\ref{sec:future-work}.

% [PENDING] Discussion of detection performance once Phase 2 is complete.
% [PENDING] Comparison with the Study 1 result, and with the reviewed
% literature, particularly Yousuf et al. (2024) on assumed thresholds and
% Boncea and Bacivarov (2016) on Prometheus-based monitoring architecture.

% =====================================================================
\subsection{Answer to RQ2}\label{sec:rq2-answer}

% [PENDING] Complete after Phases 2 and 3.
%
% Follow the structure of Section 4.6:
%   1. Restate RQ2.
%   2. State what the results show, sub-question by sub-question.
%   3. Identify the mechanisms responsible for that outcome.
%   4. State the limits of the claim, deferring the detail to Threats to
%      Results.
